\begin{tikzpicture}[
  font=\scriptsize,
  execute at begin node={\hyphenpenalty=10000\exhyphenpenalty=10000},
  lapis/.style={draw, rounded corners, fill=blue!8, text width=2.2cm, minimum height=1.05cm, align=center},
  tekan/.style={lapis, fill=orange!12},
  atas/.style={lapis, fill=green!10},
  isi/.style={draw, rounded corners, text width=7.4cm, minimum height=1.05cm, align=left},
  masalah/.style={draw, rounded corners, fill=red!22, text width=10.1cm, minimum height=0.6cm, align=center},
  ket/.style={font=\tiny, align=center},
  panah/.style={-{Stealth[length=2mm]}, thick}
]
  % Lapis disusun dari bawah supaya pembaca melihat lapis teknis sebagai fondasi.
  \node[lapis] (sintaks) at (0,0) {\textbf{1. Sintaks}};
  \node[tekan] (semantik) at (0,1.25) {\textbf{2. Semantik}};
  \node[lapis] (proses) at (0,2.5) {\textbf{3. Proses}};
  \node[atas] (tanggung) at (0,3.75) {\textbf{4. Tanggung jawab}};

  \node[isi, anchor=west] at (1.5,0) {Bentuk pesan dapat dibaca mesin.\\{\tiny\mdseries Contoh: JSON sah, elemen wajib terisi}};
  \node[isi, anchor=west] at (1.5,1.25) {Kode dan satuan bermakna sama bagi kedua pihak.\\{\tiny\mdseries Contoh: diagnosis memakai ICD-10, bukan kode lokal}};
  \node[isi, anchor=west] at (1.5,2.5) {Urutan, waktu kirim, dan cara pulih disepakati.\\{\tiny\mdseries Contoh: Encounter dikirim sebelum data klinisnya}};
  \node[isi, anchor=west] at (1.5,3.75) {Pemilik data menjawab bila isinya salah atau tidak terkirim.\\{\tiny\mdseries Contoh: Manajer Rekam Medis menandatangani rekonsiliasi}};

  \draw[panah] (-1.55,-0.5) -- node[ket, rotate=90, above] {makin sulit diperiksa mesin} (-1.55,4.25);
  \node[masalah] at (3.75,-1.05) {Status terkirim dan kode 200 hanya membuktikan lapis 1};
\end{tikzpicture}
